\subsection{The actual polarized block and the trace-one theta diagonal}

\begin{proposition}[DP47: the hyperbolic Poincar\'e pullback]\label{DP47}
Let $G$ be the actual Honda line of height $h\ge2$, or an arbitrary-unit line model under the hypotheses of DP46. Put $H=G^D$, and let
$B=\exp b$ be its invariant primitive Poincar\'e splitting in the integral frame of DP37--DP41 or DP46. Set $A=G\times H$. Suppose the integral symmetric biextension on $A\times A$ is, with axial rigidifications, the actual hyperbolic evaluation biextension
\[
 \mathcal P_{\mathrm{hyp}}((x,y),(x',y'))
       =\mathcal P(x,y')\otimes\mathcal P(x',y).
\]
Let the commutative CM action preserve this decomposition with its adjoint involution, and let $w_0+w_0^*=1$. In the block coordinates write
$w_0=(a,(1-a)^D)$, where $a$ is an integral endomorphism of $G$ commuting with its Teichm\"uller action.

Under the source convention $\mathrm{ph}=\exp(-\delta h)$, the product frame has canonical gauge
$h_0((x,y),(x',y'))=-b(x,y')-b(x',y)$.
There is an integral normalized unit $U_0$ such that
\[
 \exp(h_0(z,w_0z))=U_0(x,y)B(x,y)^{-1}.
\]
For any other source canonical gauge in any normalized integral frame,
\begin{equation}\label{eq:hyperbolicpullback}
 \exp(h(z,w_0z))=U(x,y)B(x,y)^{-1}
                       \exp Q(\log_Gx,\log_Hy),
\end{equation}
where $U$ is an integral normalized unit and $Q$ is a homogeneous quadratic form over the fraction field.
The inverse biextension convention reverses both signs coherently.
Thus for an integral section germ $f$ and integer tensor power $\kappa\ge0$,
\[
 f\exp(\kappa h_L)=\widetilde f\,B^{-\kappa}\exp(\kappa Q),
 \qquad \widetilde f=fU^\kappa\ \hbox{integral}.
\]
The primitive coefficient is exactly $-\kappa$; for DP42's unit case one requires $p\nmid\kappa$.
\end{proposition}

\begin{proof}
Write the integral addition factors as
\[
 C_G(x,x';y)=\frac{B(x+_Gx',y)}{B(x,y)B(x',y)},\qquad
 C_H(x;y,y')=\frac{B(x,y+_Hy')}{B(x,y)B(x,y')}.
\]
The actual Poincar\'e adjunction
$(a\times1)^*\mathcal P\simeq(1\times a^D)^*\mathcal P$
is integral and preserves the laws and rigidifications. In the integral frame its scalar is a normalized integral unit.
The ratio of the two rational bilinear sections after this adjunction is a rational bihomomorphism. Its logarithm is bilinear in additive logarithm coordinates: additivity in each variable forces degree one in each by comparing binomial coefficients in characteristic zero.

That bilinear ambiguity vanishes here. The adjunction is natural, $a$ commutes with the Teichm\"uller action, and the chosen frame and splitting are invariant under
$(x,y_i)\mapsto(\zeta x,\zeta^{-p^i}y_i)$.
No additive monomial $XY_i$ has trivial character, since
$0<p^i-1<p^h-1$ for $1\le i<h$.
Therefore
\[
 J_a(x,y):=\frac{B(ax,y)}{B(x,a^Dy)}
\]
is exactly an integral normalized unit, with no additional bilinear exponential.
This is the scalar of an actual integral adjunction, not an appeal to arbitrary integral-unit powers.

Put $c=1-a$. Cartier duality is additive, so
$a^Dy+_Hc^Dy=y$. Consequently
\[
 B(x,c^Dy)B(ax,y)
 =J_a(x,y)\,C_H(x;a^Dy,c^Dy)^{-1}B(x,y).
\]
All substituted addition factors are integral units. Taking inverses proves the first formula with $U_0=C_H/J_a$.
The trace-one relation therefore gives one primitive factor, with no difference of eigenvalues.
The same adjunction units verify the source's adjoint identities for $h_0$; its symmetry is literal, its axes vanish, and its addition identities are the negative coboundaries displayed above. This verifies that it is a canonical gauge in the stated convention.

In a fixed integral frame, divide the exponential addition identities for $h$ and $h_0$. Injectivity of formal exponential over the characteristic-zero field makes their difference additive in both arguments, hence bilinear in logarithm coordinates.
Symmetry and adjoint balance follow from the other gauge identities.
Every formal endomorphism becomes its tangent linear map under logarithms. The diagonal difference is therefore $L(Z,dw_0Z)$, a homogeneous quadratic form.
A change of normalized integral frame contributes the logarithm of an integral normalized unit; after integral pullback its exponential is another such unit. Adjust the frames first, then apply the bilinear-difference argument. This proves \eqref{eq:hyperbolicpullback}. Taking the ordinary integer power $\kappa$ preserves integrality of the unit and proves the last assertion.
\end{proof}

The hyperbolic hypothesis itself is geometric under a precise integral criterion. Suppose a polarized CM formal block has conjugate integral idempotents $e,e^*$ with
$A=eA\oplus e^*A=G\oplus H$, its polarization is prime to $p$ on the block, and its adjoint is CM conjugation.
For $z,z'\in eA$, balance gives
$\mathcal B(z,z')\simeq\mathcal B(z,e^*z')=\mathcal B(z,0)$,
canonically trivial; the same holds on $H\times H$.
The perfect polarization identifies the cross factor $H$ with the actual $G^D$.
Symmetry identifies the opposite cross restriction with its transpose, and biadditivity identifies the entire block with the product of these two cross restrictions. This proves the hyperbolic hypothesis integrally. It does not follow just from a slope count. The alternating finite-torsion Weil pairing does not change this symmetric-biextension transposition convention.

The exact source sign and diagonal are
\path{CocycleRightZeroThd.lean}, \path{IsCanonGauge}, lines 174--183,
\path{hL}, line 185, and \path{nabVal}, lines 187--189.
\path{WickVanishing/ExpConjConnectionIntegral.lean}, line 40, explicitly assumes $w_0+\operatorname{star}(w_0)=1$.
\path{Development/AbelianVarietyTheory/Proofs/FormalGroupCharts/CanonGaugeSubIsBiadd.lean}, lines 16--60, supplies the gauge-difference identities; \path{CanonGaugeIndependenceAll.lean}, lines 19--26, supplies the specified selected-value comparison.
The actual pairing and its functoriality are Zink (216)--(220), Proposition~121, (243) and Theorem~122 \cite{Zink}.

On a primitive mixed logarithmic line, at odd $p$, $h>2$ and $p\nmid\kappa$, the resulting expression is
$F(t)\exp(-\kappa\beta(ut,\mathbf vt)+ct^2)$.
The necessary condition is $F\ne0$ as a restricted series; global nonvanishing alone does not prove this.
DP48 treats the quadratic term. These results identify the actual local kernel under the polarized-block hypotheses, while the prescribed global path, selected harmonic module, nonzero minors and cost comparison remain separate inputs.

